\apendice{Documentación técnica de programación}

\section{Introducción}

En este apartado se describe la estructura técnica del sistema desarrollado, así como los aspectos necesarios para su instalación, ejecución y mantenimiento.

El sistema está compuesto por dos addons diferenciados pero complementarios:

\begin{itemize}
    \item Un addon de recopilación de datos, encargado de registrar la actividad del usuario durante el proceso de modelado.
    \item Un addon de análisis y visualización, encargado de procesar la información generada y representar métricas mediante gráficos y visualizaciones tridimensionales.
\end{itemize}

Ambos componentes trabajan conjuntamente mediante el uso de archivos CSV como mecanismo de intercambio de información.

En este anexo se describen la organización interna del proyecto, la estructura de directorios, los módulos principales de cada addon y los procedimientos necesarios para su instalación y ejecución.

\section{Estructura de directorios}

El proyecto se divide en dos componentes principales correspondientes a los addons desarrollados.

\subsection{Addon de recopilación de datos}

El addon de recopilación se implementa principalmente como un script Python ejecutable dentro de Blender.

Su estructura principal es la siguiente:

\begin{itemize}
    \item \textbf{DeteccionDatosV4.py}: archivo principal encargado de:
    \begin{itemize}
        \item Registrar handlers de Blender.
        \item Detectar eventos en la escena.
        \item Capturar propiedades geométricas y espaciales.
        \item Generar archivos CSV.
        \item Gestionar controles Start/Stop y paneles de depuración.
    \end{itemize}
\end{itemize}

Este enfoque facilita su ejecución directa desde el editor de scripting de Blender.

\subsection{Addon de análisis y visualización}

El segundo addon se distribuye como un paquete modular compuesto por distintos archivos Python.

La estructura principal del addon es la siguiente:

\begin{itemize}

    \item \textbf{\_\_init\_\_.py}: archivo de inicialización e integración del addon con Blender.

    \item \textbf{metrics\_ui.py}: definición de paneles, interfaz gráfica y controles de visualización de métricas.

    \item \textbf{graphs.py}: generación de gráficos y representaciones visuales.

    \item \textbf{utils.py}: funciones auxiliares para lectura y procesamiento de datos.

    \item \textbf{operators.py}: definición de operadores personalizados utilizados por Blender.

    \item \textbf{visualization.py}: generación de geometría tridimensional y representaciones visuales dentro de la escena.

    \item \textbf{site-packages/}: dependencias externas necesarias para el funcionamiento del addon.

\end{itemize}

La estructura modular facilita el mantenimiento, la reutilización de código y la ampliación futura del sistema.

\section{Manual del programador}

El sistema se ha desarrollado siguiendo un enfoque modular y orientado a eventos.

\subsection{Addon de recopilación}

El addon de recopilación utiliza handlers de Blender para detectar automáticamente cambios dentro de la escena.

Su funcionamiento principal consiste en:

\begin{itemize}
    \item Detectar cambios mediante eventos del \textit{dependency graph}.
    \item Identificar el tipo de operación realizada.
    \item Capturar propiedades relevantes del objeto afectado.
    \item Registrar la información en archivos CSV.
\end{itemize}

Entre las funcionalidades implementadas destacan:

\begin{itemize}
    \item Detección de creación y eliminación de objetos.
    \item Detección de cambios geométricos en \textit{Edit Mode}.
    \item Captura de coordenadas UV.
    \item Detección de operadores como duplicado, pegado o merge.
    \item Control de ejecución mediante Start/Stop.
    \item Logging basado en cambios.
\end{itemize}

El sistema implementa mecanismos para evitar registros redundantes, almacenando únicamente cambios significativos en el estado de la escena.

\subsection{Addon de análisis y visualización}

El addon de análisis se organiza de forma modular, separando claramente procesamiento, interfaz y representación visual.

Sus principales componentes funcionales son:

\begin{itemize}

    \item \textbf{Carga de datos}: lectura y validación de archivos CSV generados por el addon de recopilación.

    \item \textbf{Procesamiento de métricas}: cálculo de estadísticas temporales, espaciales y estructurales.

    \item \textbf{Visualización 2D}: generación de gráficos estadísticos mediante librerías externas como \textit{matplotlib}.

    \item \textbf{Visualización 3D}: creación dinámica de geometría y representaciones analíticas dentro de Blender.

    \item \textbf{Interfaz de usuario}: paneles y operadores personalizados integrados en la interfaz de Blender.

\end{itemize}

El addon permite representar métricas directamente dentro de la escena tridimensional, facilitando la interpretación visual de la información recopilada.

\section{Compilación, instalación y ejecución del proyecto}

\subsection{Addon de recopilación}

El addon de recopilación se ejecuta directamente desde Blender siguiendo estos pasos:

\begin{enumerate}
    \item Abrir Blender.
    \item Acceder al espacio de trabajo \textit{Scripting}.
    \item Cargar el archivo \texttt{DeteccionDatosV4.py}.
    \item Ejecutar el script.
    \item Registrar el addon mediante la opción correspondiente.
    \item Guardar el archivo \texttt{.blend}.
\end{enumerate}

Una vez ejecutado, el sistema solicita el consentimiento del usuario para la recopilación de datos. A partir de ese momento, la captura de eventos se realiza automáticamente durante el uso de Blender.

\subsection{Addon de análisis y visualización}

El addon de análisis se instala como un addon estándar de Blender.

El procedimiento de instalación es el siguiente:

\begin{enumerate}
    \item Comprimir el addon en formato \texttt{.zip}.
    \item Acceder a \textit{Edit > Preferences > Add-ons}.
    \item Seleccionar la opción \textit{Install}.
    \item Importar el archivo \texttt{.zip}.
    \item Activar el addon.
\end{enumerate}

Una vez instalado, el usuario puede acceder al panel de análisis y cargar los archivos CSV generados previamente.

\section{Dependencias externas}

El addon de análisis incorpora varias librerías externas utilizadas para el procesamiento y visualización de datos.

Entre las principales dependencias utilizadas destacan:

\begin{itemize}
    \item \textbf{matplotlib}: generación de gráficos estadísticos.
    \item \textbf{numpy}: operaciones numéricas y normalización de métricas.
    \item \textbf{csv}: lectura y escritura de archivos CSV.
\end{itemize}

Estas dependencias se integran dentro del addon mediante el directorio \texttt{site-packages}.

\section{Pruebas del sistema}

Se realizaron diferentes pruebas para validar el correcto funcionamiento de ambos addons.

\subsection{Pruebas del addon de recopilación}

\begin{itemize}
    \item Registro correcto de creación y eliminación de objetos.
    \item Detección de modificaciones en \textit{Edit Mode}.
    \item Captura de operaciones como Shift+D, Ctrl+V y Merge.
    \item Validación de coordenadas UV.
    \item Comprobación del funcionamiento de controles Start/Stop.
    \item Verificación de generación correcta de archivos CSV.
    \item Validación del sistema de logging basado en cambios.
\end{itemize}

\subsection{Pruebas del addon de análisis}

\begin{itemize}
    \item Lectura correcta de archivos CSV.
    \item Verificación del cálculo de métricas.
    \item Validación de gráficos estadísticos.
    \item Generación correcta de visualizaciones 3D.
    \item Comprobación de integración de objetos visuales en Blender.
    \item Validación de normalización y representación de métricas.
    \item Verificación de funcionamiento de paneles y operadores personalizados.
\end{itemize}

Estas pruebas permitieron garantizar el correcto funcionamiento del sistema completo, tanto en la captura de datos como en su posterior análisis y representación visual.

\section{Mantenimiento y ampliación}

La arquitectura modular adoptada facilita el mantenimiento y la ampliación futura del sistema.

La separación entre captura, procesamiento y visualización permite incorporar nuevas funcionalidades sin modificar significativamente el resto de componentes.

Entre las posibles ampliaciones futuras destacan:

\begin{itemize}
    \item Incorporación de nuevas métricas de análisis.
    \item Soporte para otros formatos de almacenamiento.
    \item Exportación automática de resultados.
    \item Integración con sistemas de analíticas de aprendizaje.
    \item Mejora de las visualizaciones tridimensionales.
\end{itemize}

El diseño modular del sistema favorece su reutilización y evolución en futuras investigaciones o desarrollos relacionados con el análisis de interacción en entornos tridimensionales.